\section{Diagnostic Audit}
\label{app:diagnostic_audit}

The diagnostic audit examines whether empirical evidence matches the process-level claim made by a paper. It does not estimate prevalence in the entire literature. The sample is purposively stratified to cover the four functional categories, major subcategories, older and recent methods, fixed-corpus and open-world settings, and both component-level and end-to-end claims.

\subsection{Sampling}
The \AuditSize-paper selection contains 15 initial-retrieval papers, 25 follow-up-retrieval papers, 20 evidence-integration papers, and 10 feedback-learning papers. Follow-up retrieval receives the largest allocation because it contains the widest range of trajectory-control mechanisms and because final-answer accuracy is particularly insufficient for distinguishing their effects. The audit list was frozen before quantitative summaries were produced.

\subsection{Coding Fields}
Table~\ref{tab:audit_codebook} summarizes the core codebook. Each field is coded from the full paper, including appendices where relevant.

\begin{table*}[tbp]
\centering
\caption{Diagnostic-audit codebook.}
\label{tab:audit_codebook}
\small
\begin{tabularx}{\textwidth}{@{}p{3.0cm}p{3.6cm}Y@{}}
\toprule
Field & Allowed or typical values & Coding question \\
\midrule
Primary claim & retrieval; search control; integration; grounding; efficiency; cross-task learning & What component does the paper claim to improve? \\
Task and evidence setting & fixed corpus; graph; table/text; multimodal; open web; repeated tasks & Under what information environment is the claim tested? \\
End-to-end outcome & EM/F1; accuracy; report score; human preference; task reward & Is the final task outcome reported? \\
Retrieval diagnostics & Recall@$k$; MRR; nDCG; bridge recall; complete-chain recall & Does the paper directly measure evidence acquisition? \\
Trajectory diagnostics & subquestion quality; hop success; action validity; stopping; recovery; trace release & Does the paper measure the behavior it uses to control later search? \\
Integration diagnostics & evidence membership; compression fidelity; ordering; citation; evidence intervention & Does the paper isolate how acquired evidence is transformed and used? \\
Learning diagnostics & learning curve; transfer; memory or skill retrieval; forgetting; negative transfer & Does the paper demonstrate improvement across tasks rather than only within one episode? \\
Efficiency & retrieval calls; model calls; tokens; latency; wall-clock time; monetary cost & Is the additional computation needed for the gain reported? \\
Controls and interventions & fixed candidate pool; fixed evidence set; reader control; evidence removal; component ablation & Can the reported gain be attributed to the claimed component? \\
Evidence locator & page, section, table, figure, appendix, or quotation paraphrase & Where in the source is the coding supported? \\
Unresolved gap & free-text diagnosis & What evidence would be needed to support the claim more directly? \\
\bottomrule
\end{tabularx}
\end{table*}

\subsection{Coding Values}
Binary-looking fields use four principal values: \emph{reported}, \emph{not reported}, \emph{not applicable}, and \emph{unclear}. We additionally use \emph{partial} when a paper reports a closely related diagnostic but does not fully expose the claimed process. These values are not collapsed. ``Unclear'' is used when the paper's terminology or experimental setup does not permit a defensible decision. ``Not applicable'' is reserved for cases in which the diagnostic is genuinely irrelevant to the stated claim, not merely absent.

A positive code requires a source locator. For example, a method is credited with trajectory evaluation only when the paper reports a query-, action-, hop-, or stopping-level result; showing one qualitative example is recorded separately from aggregate evaluation. Likewise, a paper is credited with causal evidence-use analysis only when evidence is removed, replaced, reordered, perturbed, or otherwise intervened on while relevant controls are held fixed.

\subsection{Claim-Aligned Interpretation}
The audit does not require every paper to report every field. A dense retriever may not need cross-task learning curves, and a skill-learning method may not need MRR. The requirement is alignment: a claim about better retrieval should include retrieval diagnostics; a claim about better long-horizon search should expose trajectories and costs; a claim about evidence use should include membership or intervention evidence; and a claim about learning from experience should compare future-task performance against a matched no-experience baseline.

\subsection{Completed Calibration Batch}
The first completed batch contains eight papers, with two selected from each functional category. Table~\ref{tab:audit_completed_batch} records the source-located interpretation used to calibrate the codebook. The corresponding machine-readable records are released in \texttt{audit/audit\_batch\_01\_completed.csv}. These eight cases are not a random sample and are not used to compute literature-wide reporting percentages.

\begin{table*}[tbp]
\centering
\caption{First completed full-text audit batch. ``Partial'' indicates that the paper reports a related diagnostic without fully exposing the target process.}
\label{tab:audit_completed_batch}
\scriptsize
\begin{tabularx}{\textwidth}{@{}p{2.15cm}p{2.25cm}p{2.10cm}Y Y@{}}
\toprule
Paper & Primary claim & Direct diagnostics & Source locator & Unresolved diagnostic \\
\midrule
DPR~\cite{karpukhin2020dpr} & Retrieval effectiveness and efficiency & Retrieval and end-to-end outcomes; efficiency; controls & pp.~4, 6--7; Tables~2, 4 & Bridge and complete-chain recall are outside the original single-hop task and should not be inferred \\
ColBERT~\cite{khattab2020colbert} & Ranking effectiveness and efficiency & Retrieval and ranking metrics; latency; FLOPs; fixed-pool comparison & pp.~6--8; Tables~1--2 & No downstream QA or multi-hop support-chain evaluation for the general ranking claim \\
MDR~\cite{xiong2021mdr} & Multi-hop passage-sequence retrieval & Support-passage EM, answer recall, R@$k$, downstream outcome, component ablations & pp.~3--5; Tables~2--3 & Recovery, stopping, and per-instance query/action traces remain partial; some speed comparisons are estimated \\
IRCoT~\cite{trivedi2023ircot} & Reasoning-guided iterative retrieval & Alternating steps, explicit stopping rule, fixed paragraph cap, retrieval recall, answer F1, factual-error analysis & pp.~3--6; Fig.~2; Figs.~3--4; Appendix~C & No aggregate invalid-query, recovery, premature-stopping, latency, or monetary-cost evaluation \\
LLMLingua~\cite{jiang2023llmlingua} & Prompt compression with task preservation & Compression ratio, token count, task outcome, latency, and component ablations & pp.~5--8; Tables~2, 6 & Relation support and provenance preservation are inferred indirectly rather than measured by intervention \\
Lost in the Middle~\cite{liu2024lostmiddle} & Position sensitivity in long-context use & Fixed evidence membership with controlled changes to position and context length & pp.~1--8; Figs.~1, 5, 7--10 & Does not test multi-hop relation preservation or claim-level source grounding \\
Reflexion~\cite{shinn2023reflexion} & Verbal feedback and episodic memory across retries & Repeated-trial learning curves and memory/reflection ablations & pp.~3--6; Figs.~2, 4; Appendix~D.4 & Held-out cross-task transfer, negative transfer, forgetting, and matched context overhead are not established \\
ExpeL~\cite{zhao2024expel} & Cross-task experience extraction and recall & Unseen-task reuse, same-distribution gains, one forward-transfer setting, and component ablations & pp.~2--9, 35; Tables~1--3, 6 & Negative transfer, forgetting, stale memory, and long-term library growth are not tested \\
\bottomrule
\end{tabularx}
\end{table*}

The calibration produced three codebook decisions. First, an irrelevant field is coded \emph{not applicable}, not \emph{not reported}; DPR and ColBERT are not penalized for omitting cross-task or multi-hop trajectory measures that are outside their original claims. Second, a method can receive a \emph{partial} trajectory code when it defines hop transitions or stopping but does not aggregate failure, recovery, or action-cost statistics. Third, downstream task preservation is not equivalent to causal preservation of evidence relations or provenance. This distinction separates LLMLingua's strong compression and latency evaluation from the controlled position intervention used by Lost in the Middle.

\subsection{Current Status and Reliability Boundary}
Eight of the \AuditSize{} selected papers have completed source-located coding; 62 remain in the frozen selection. The manuscript therefore continues to describe the resource as a \emph{70-paper audit selection} and labels the completed subset explicitly. No category-level reporting percentage will be added until the relevant claim-aligned denominator has been fully coded.

The current coding is single-coder work. Before publication, a stratified subset should be independently double-coded, disagreements should be adjudicated, and the agreement statistic should be reported together with the sample and field definitions. No reliability statistic is inferred from agreement on metadata or from automated checks.

\subsection{Planned Quantitative Summaries}
After full-text coding, the audit will report category-wise coverage of retrieval, trajectory, integration, learning, efficiency, and intervention diagnostics. Counts will use the relevant claim-aligned denominator rather than all 70 papers indiscriminately. Qualitative synthesis will focus on recurrent gaps, such as answer-only evaluation for search methods, missing fixed-budget controls, absent source-use interventions, and cross-task learning claims evaluated only on repeated or closely matched tasks.
